\subsection{Compactness of the smeared plaquette field from the free energy}
\label{rU7:sec}

This subsection proves an unconditional compactness statement for the
renormalised, smeared plaquette field along an arbitrary sequence of
couplings and spacings, and identifies exactly which property of the
partition function controls it. The tool is the chessboard estimate\footnote{Reflection positivity in the planes between lattice sites gives a Cauchy--Schwarz inequality for each reflection; applying it repeatedly in every direction bounds the expectation of a product of translates of one local observable by the product of the geometric means of the same observable spread over the whole torus~\cite{FILS78}.} of the
Wilson measure on a torus, which reduces every exponential moment of a
smeared plaquette field to the free energy of the same torus at a shifted
coupling. Nothing below assumes the averaged condition or any decay, and
everything is proved for every $\beta_k\ge0$, except where the
short-distance law (AF$^\square$) of Sec.~\ref{rU7:af} or the stability
hypothesis~(B) of Sec.~\ref{rU7:balaban} is stated as a hypothesis. Multiple reflections were used
for the same purpose by O'Carroll and Faria da Veiga~\cite{OCFdV20}, who bound
the generating functional of a fixed number of unsmeared gauge-invariant
plaquette fields uniformly in the spacing and the volume; the bound below is for fields smeared over
$\asymp a^{-4}$ plaquettes, and it is an identity for the constant smearing,
which is what makes it a criterion rather than an estimate. The relation to the inputs of
Sec.~\ref{ssec:rEmissing} is in Sec.~\ref{rU7:af} and the relation to the
renormalisation group bounds of Ba{\l}aban in Sec.~\ref{rU7:balaban}.

\emph{Setting.} Let $G\subset\U(N)$ be a closed subgroup, $G\ne\{1\}$, with
the Wilson weight $\exp(N\beta\sum_p\ReTr Q_p)$ of Sec.~\ref{sec:intro},
$\beta\ge0$.
Fix $L>0$ and let $\mathbb T_L=(\R/L\Z)^4$. For an even integer $M\ge4$ let
$\Lambda_M=(\Z/M\Z)^4$, $a=L/M$, $\mu_{M,\beta}$ the Wilson measure of
$\Lambda_M$ at $\beta$, $S=\sum_p\ReTr Q_p$ the sum over the $6M^4$
plaquettes, and
\[
Z_M(\beta)=\int e^{N\beta S}\,dQ,\qquad p_M(\beta)=M^{-4}\log Z_M(\beta),
\]
with normalised Haar measure $dQ$ on $G^{E(\Lambda_M)}$; $Z_M(\beta)$ and
$p_M(\beta)$ are defined by these formulas for every real $\beta$, and
$\mu_{M,\beta}$ is the probability measure with density $e^{N\beta S}/Z_M(\beta)$.
The function $p_M$ is real analytic and convex on $\R$, with
\begin{align*}
p_M'(\beta)&=\frac N{M^4}\,\mu_{M,\beta}(S)=6N\,\mu_{M,\beta}(\ReTr Q_p),\\
p_M''(\beta)&=\frac{N^2}{M^4}\Var_{\mu_{M,\beta}}(S)>0,
\end{align*}
(the second equality in the first line by the invariance of $\mu_{M,\beta}$
under translations and the hypercubic group, which make the six
orientations equivalent), and $|p_M'|\le6N^2$ because $|\ReTr Q_p|\le N$; the
variance is positive
because $S$ is not constant on $G^{E(\Lambda_M)}$ ($\ReTr g<N$ for $g\in
G\setminus\{1\}$, so changing one link of the identity configuration to $g$
lowers $S$, and $S$ is continuous while Haar measure charges every open
set). Its second-order remainder
\begin{align*}
R_M(\beta;h)&=p_M(\beta+h)-p_M(\beta)-h\,p_M'(\beta)\\
&=\int_0^h\bigl(p_M'(\beta+s)-p_M'(\beta)\bigr)\,ds\ \ge0
\end{align*}
is convex in $h$, vanishes at $h=0$, is strictly increasing on $h\ge0$ and
strictly decreasing on $h\le0$, and satisfies $R_M(\beta;sh)\le sR_M(\beta;h)$
for $s\in[0,1]$.

For a site $x\in\Lambda_M$ let $C_x=x+[0,1]^4$ be the closed unit cube with
lower corner $x$, $m_x=x+\frac12\mathbf 1$ its centre, and
\[
\psi_p=N-\ReTr Q_p\in[0,2N],\qquad
\Phi^\square_x=\sum_{p\subset C_x}\psi_p ,
\]
the sum over the $24$ two-dimensional faces of $C_x$; $\Phi^\square_x$ is a
gauge-invariant function of the links of $C_x$ and is invariant under the
hypercubic group of $C_x$. Every plaquette lies in exactly four unit cubes
(two choices of the lower corner in each of its two transverse directions),
so
\[
\sum_{x\in\Lambda_M}\Phi^\square_x=4\sum_p\psi_p=24NM^4-4S .
\tag{rU7.1}\label{rU7:sum}
\]
For a bounded function $f$ on $\mathbb T_L$, $Z>0$ and the measure
$\mu=\mu_{M,\beta}$ with $c=\mu(\Phi^\square_0)=24\mu(\psi_p)$, put
\begin{align*}
\mathcal A(f)&=\mathcal A_{M,\beta,Z}(f)=Z\sum_{x\in\Lambda_M}a^4f(am_x)\,
\bigl(\Phi^\square_x-c\bigr),\\
\|f\|_{1,a}&=\sum_xa^4|f(am_x)| ,
\end{align*}
so that $\|f\|_{1,a}\le L^4\|f\|_\infty$, and $\|f\|_{1,a}\to\|f\|_{L^1}$ as
$a\to0$ for Riemann integrable $f$. The $24$-plaquette star field $\Phi_x$
(the sum of $\psi_p$ over the plaquettes with a corner at $x$, the example
of a local field in Sec.~\ref{ssec:limitE}) and $\Phi^\square$ are two smearings of the same
plaquette field, $\sum_xf(ax)\Phi_x=\sum_p\psi_p\sum_{\text{corners }y\text{
of }p}f(ay)$ and $\sum_xf(am_x)\Phi^\square_x=\sum_p\psi_p\sum_{\text{cubes }
C_x\supset p}f(am_x)$, with weights that differ by at most $8\,\mathrm w_f(a)$, $\mathrm w_f$ the
modulus of continuity of $f$ (corners and centres of the four cubes are at
distance $a/\sqrt2$ from the centre of $p$); the results below are stated for
$\Phi^\square$ because its chessboard reflections are isotropic
(Lemma~\ref{rU7:chess}).

\begin{lemma}[the constant source is a shift of the coupling]\label{rU7:const}
For every $\lambda\in\R$,
\[
\mu_{M,\beta}\Bigl(\exp\Bigl(\lambda\sum_{x\in\Lambda_M}(\Phi^\square_x-c)
\Bigr)\Bigr)=\exp\Bigl(M^4R_M\bigl(\beta;-\tfrac{4\lambda}N\bigr)\Bigr) .
\]
\end{lemma}

\begin{proof}
By~\eqref{rU7:sum} and $M^4c=\mu(\sum_x\Phi^\square_x)=24NM^4-4\mu(S)$,
$\sum_x(\Phi^\square_x-c)=-4(S-\mu(S))$. Since $e^{N\beta S}e^{-4\lambda S}=
e^{N(\beta-4\lambda/N)S}$,
\begin{align*}
\mu\bigl(e^{-4\lambda(S-\mu(S))}\bigr)&=e^{4\lambda\mu(S)}\,
\frac{Z_M(\beta-4\lambda/N)}{Z_M(\beta)}\\
&=\exp\Bigl(M^4\Bigl[\tfrac{4\lambda}N\,p_M'(\beta)\\
&\qquad\qquad+p_M\bigl(\beta-\tfrac{4\lambda}N\bigr)-p_M(\beta)\Bigr]\Bigr),
\end{align*}
using $\mu(S)=M^4p_M'(\beta)/N$; the bracket is
$R_M(\beta;-4\lambda/N)$.
\end{proof}

\begin{lemma}[chessboard bound for the plaquette field]\label{rU7:chess}
Let $M\ge4$ be even, $\beta\ge0$, and $\lambda:\Lambda_M\to\R$. Then
\[
\mu_{M,\beta}\Bigl(\exp\Bigl(\sum_x\lambda_x(\Phi^\square_x-c)\Bigr)\Bigr)
\le\exp\Bigl(\sum_{x\in\Lambda_M}R_M\bigl(\beta;-\tfrac{4\lambda_x}N\bigr)
\Bigr) .
\]
\end{lemma}

\begin{proof}
\emph{Reflection positivity.} Let $\vartheta$ be the reflection of
$\Lambda_M$ in the hyperplane $\{x_\mu=c_0\}$, $c_0\in\Z$, which also fixes
the antipodal hyperplane $\{x_\mu=c_0+M/2\}$; it maps links to links,
preserving orientation for links in the two hyperplanes and otherwise
mapping a link to a link, and plaquettes to plaquettes. Let $\mathfrak A_+$
be the algebra of continuous functions of the links with both endpoints in
the closed half $\{c_0\le x_\mu\le c_0+M/2\}$, $\mathfrak A_-=\mathfrak A_+
\circ\vartheta$, and $(\vartheta F)(Q)=F(Q\circ\vartheta)$. Every plaquette
lies in one closed half or in both, the latter exactly when it lies in one
of the two hyperplanes, so $S=S_++\vartheta S_+$ with $S_+\in\mathfrak A_+$
the sum of $\ReTr Q_p$ over the plaquettes in the positive closed half,
those in the hyperplanes counted with weight $\frac12$. Haar measure is
invariant under $Q_e\mapsto Q_{\vartheta e}$ and under inversion, so,
integrating first over the links not in the hyperplanes,
\begin{align*}
\int\overline{\vartheta F}\,F\,e^{N\beta S}\,dQ&=\int dQ_0\,
\Bigl|\int F\,e^{N\beta S_+}\,dQ_+\Bigr|^2\ge0\\
&\qquad(F\in\mathfrak A_+) ,
\end{align*}
where $dQ_0$ is Haar measure on the links in the two hyperplanes and $dQ_+$
on the other links of the positive closed half,
which is the reflection positivity of $\mu_{M,\beta}$ in every hyperplane of
sites, proved by Osterwalder and Seiler~\cite{OSei78}.

\emph{Chessboard estimate.} The torus is tiled by the closed unit cubes
$C_x$, neighbouring cubes share a face, and the reflections in the
hyperplanes containing the faces are those of the preceding paragraph; a
function of the links of $C_x$ lies in $\mathfrak A_+$ for every such
hyperplane with $C_x$ on its positive side. For such block functions $F_x$
the chessboard estimate of Fr\"ohlich, Israel, Lieb and
Simon~\cite{FILS78}, in the form of~\cite[Theorem~5.8]{Bis09} with blocks of
side $B=1$ and $M/B=M$ even (reflections through planes of sites), gives
\[
\mu\Bigl(\prod_xF_x\Bigr)\le\prod_x\mu\Bigl(\prod_y\vartheta_{y\leftarrow x}
F_x\Bigr)^{1/M^4},
\]
where $\vartheta_{y\leftarrow x}F_x$ is $F_x$ carried into the cube $C_y$ by
a chain of reflections in faces (independent of the chain, since
reflections in different directions commute and $\vartheta^2=1$).
In~\cite{Bis09} the variables sit on sites and a $B$-block function depends
on the sites of $\{0,\dots,B\}^d$; here they sit on links and a block
function depends on the links with both endpoints in a closed unit cube, and
the proof of~\cite[Theorem~5.8]{Bis09} carries over word for word. It uses
only (a) the Cauchy--Schwarz inequality for the positive semidefinite
symmetric form $(F,F')\mapsto\mu(F\,\vartheta F')$ on real elements of
$\mathfrak A_+$, which is the case $F$ real of the display above together
with the invariance of $\mu$ under $\vartheta$; (b) the algebraic fact that
$\vartheta$ maps $\mathfrak A_+$ onto $\mathfrak A_-$ as an algebra
homomorphism and that products of block functions in one closed half lie in
$\mathfrak A_+$, the links in the hyperplane belonging to both halves; and
(c) the cyclic invariance of the multilinear functional in its proof, which
follows from the invariance of $\mu$ under the reflections of the torus.
All three hold for $\mu_{M,\beta}$, which is invariant under every
reflection and translation of $\Lambda_M$ because Haar measure is invariant
under permutations and inversions of the link variables and $S$ is invariant
under these maps. Take $F_x=\exp(\lambda_x\Phi^\square_x)>0$. A reflection in a face
hyperplane maps the $24$ faces of $C_x$ onto the $24$ faces of the image
cube and preserves $\ReTr Q_p$, since $\ReTr Q^{-1}=\ReTr Q$ for unitary $Q$;
so $\vartheta_{y\leftarrow x}F_x=\exp(\lambda_x\Phi^\square_y)$ and
$\prod_y\vartheta_{y\leftarrow x}F_x=\exp(\lambda_x\sum_y\Phi^\square_y)$.
Multiplying both sides of the chessboard estimate by the constant
$\exp(-\sum_x\lambda_xc)=\prod_x\exp(-\lambda_xM^4c)^{1/M^4}$ and applying
Lemma~\ref{rU7:const} to each factor gives the claim.
\end{proof}

For $Z>0$ and $T\ge0$ put
\[
\Xi_{M,\beta,Z}(T)=a^{-4}\max\Bigl\{R_M\bigl(\beta;\tfrac{4TZa^4}N\bigr),
\ R_M\bigl(\beta;-\tfrac{4TZa^4}N\bigr)\Bigr\},
\]
the larger second-order remainder of the free energy, per unit physical
volume, at a coupling shift of size $4TZa^4/N$.

\begin{theorem}[exponential moments and the free energy]\label{rU7:torus}
Let $M\ge4$ be even, $\beta\ge0$, $Z>0$, and $\mu=\mu_{M,\beta}$.

\textup{(i)} For every bounded $f\ne0$ on $\mathbb T_L$ and $t\in\R$,
\[
\log\mu\bigl(e^{t\mathcal A(f)}\bigr)\le\frac{\|f\|_{1,a}}{\|f\|_\infty}\,
\Xi_{M,\beta,Z}\bigl(|t|\,\|f\|_\infty\bigr) .
\]

\textup{(ii)} For $f\equiv1$ there is equality up to the choice of sign:
\[
\log\mu\bigl(e^{\pm t\mathcal A(1)}\bigr)=L^4a^{-4}R_M\bigl(\beta;\mp
\tfrac{4tZa^4}N\bigr),
\]
so $\max_\pm\log\mu(e^{\pm t\mathcal A(1)})=L^4\,\Xi_{M,\beta,Z}(|t|)$.

\textup{(iii)} For every bounded $f$,
\[
\begin{split}
\mu\bigl(\mathcal A(f)^2\bigr)&\le\frac{16Z^2a^4p_M''(\beta)}{N^2}\,
\|f\|_{1,a}\|f\|_\infty\\
&=\frac{\|f\|_{1,a}\|f\|_\infty}{L^4}\,\mu\bigl(\mathcal A(1)^2\bigr).
\end{split}
\]

\textup{(iv)} Consequently, for a sequence $(M_k,\beta_k,Z_k)$ with $M_k$
even, $a_k=L/M_k$ and $\mu_k=\mu_{M_k,\beta_k}$, $\mathcal A_k=\mathcal
A_{M_k,\beta_k,Z_k}$, the following are equivalent:
\begin{enumerate}[label=\textup{(\alph*)},leftmargin=*]
\item $\sup_k\mu_k(e^{t\mathcal A_k(1)})<\infty$ for every $t\in\R$;
\item $\sup_k\Xi_{M_k,\beta_k,Z_k}(T)<\infty$ for every $T>0$;
\item $\sup_k\mu_k(e^{t\mathcal A_k(f)})<\infty$ for every $t\in\R$ and every
bounded $f$ on $\mathbb T_L$.
\end{enumerate}
When they hold, for every $n\ge1$, $T>0$ and bounded $f\ne0$,
\[
\mu_k\bigl(|\mathcal A_k(f)|^n\bigr)\le2\,n!\,T^{-n}\exp\Bigl(\frac{\|f\|_{1,
a_k}}{\|f\|_\infty}\sup_k\Xi_{M_k,\beta_k,Z_k}(T\|f\|_\infty)\Bigr).
\]
\end{theorem}

\begin{proof}
(i) Apply Lemma~\ref{rU7:chess} with $\lambda_x=tZa^4f(am_x)$ and put
$h_x=-4\lambda_x/N$, $H=4|t|Za^4\|f\|_\infty/N$, so $|h_x|\le H$. By
convexity and $R_M(\beta;0)=0$, $R_M(\beta;h_x)\le(|h_x|/H)R_M(\beta;
\operatorname{sgn}(h_x)H)\le(|f(am_x)|/\|f\|_\infty)\max_\pm R_M(\beta;\pm H)$.
Summing over $x$,
\begin{align*}
\sum_xR_M(\beta;h_x)&\le\frac{\sum_x|f(am_x)|}{\|f\|_\infty}\max_\pm R_M(\beta;
\pm H)\\
&=\frac{\|f\|_{1,a}}{a^4\|f\|_\infty}\max_\pm R_M(\beta;\pm H),
\end{align*}
which is the right side of (i), since $a^{-4}\max_\pm R_M(\beta;\pm H)=
\Xi_{M,\beta,Z}(|t|\|f\|_\infty)$.

(ii) is Lemma~\ref{rU7:const} with $\lambda=\pm tZa^4$, as $M^4=L^4a^{-4}$.

(iii) Since $\mathcal A(f)$ is a bounded centred random variable, the left
side of (i) is $\frac12t^2\mu(\mathcal A(f)^2)+O(|t|^3)$; since $p_M$ is
smooth, $R_M(\beta;\pm H)=\frac12H^2p_M''(\beta)+O(H^3)$, and $H=4|t|Za^4
\|f\|_\infty/N$, so the right side is $\frac12(\|f\|_{1,a}/\|f\|_\infty)a^{-4}
H^2p_M''(\beta)+O(|t|^3)$. Dividing (i) by $t^2/2$ and letting $t\to0$ gives
$\mu(\mathcal A(f)^2)\le(\|f\|_{1,a}/\|f\|_\infty)\,a^{-4}\,(4Za^4\|f\|_\infty
/N)^2p_M''(\beta)$, the first bound. By~\eqref{rU7:sum}, $\mathcal A(1)=
-4Za^4(S-\mu(S))$, so $\mu(\mathcal A(1)^2)=16Z^2a^8\Var(S)=16Z^2a^4L^4
p_M''(\beta)/N^2$, which gives the second form.

(iv) (c)$\Rightarrow$(a) is trivial; (a)$\Rightarrow$(b) is (ii), since (a)
for $\pm t$ bounds $L^4\Xi_k(|t|)$; (b)$\Rightarrow$(c) is (i) with
$\|f\|_{1,a_k}/\|f\|_\infty\le L^4$. For the moment bound, $|x|^n/n!\le
e^{T|x|}T^{-n}\le T^{-n}(e^{Tx}+e^{-Tx})$, and (i) bounds both exponential
moments.
\end{proof}

\begin{remark}\label{rU7:torusremark}
Part (iv) is a sharp criterion: the whole family of smeared plaquette
fields on the torus of side $L$, renormalised by $Z_k$, has uniformly bounded
exponential moments if and only if the constant smearing does, if and only
if the second-order Taylor remainder of the free energy per unit physical
volume at coupling shifts of order $Z_ka_k^4$ stays bounded. Since
$M^4[R_M(\beta;h)+R_M(\beta;-h)]=\log Z_M(\beta+h)-2\log Z_M(\beta)+\log
Z_M(\beta-h)$ and both remainders are nonnegative, (b) is equivalent to
\[
\sup_kL^{-4}\Bigl|\Delta^2_{h_k(T)}\log Z_{M_k}(\beta_k)\Bigr|<\infty\quad
(T>0),
\]
with $h_k(T)=4TZ_ka_k^4/N$,
a bound on the second difference of the logarithm of the partition
function at three couplings. This is the exact form in which a statement
about partition functions controls correlations of the plaquette field
(Sec.~\ref{rU7:balaban}). Part (iii) shows that the plaquette susceptibility
$p_M''(\beta)/N^2=M^{-4}\Var(S)=\sum_{p}\sum_{p'}\Cov(\ReTr Q_p,\ReTr Q_{p'})$,
the inner sum over all plaquettes $p'$ of the torus and the outer one over
the six plaquettes $p$ with lower corner $0$ (one for each orientation), is,
up to the factor $16Z^2a^4L^4$, the variance of the constant smearing and dominates every
smeared variance: no smearing, however singular or however close to two
coinciding points, has a larger variance per unit of $\|f\|_{1,a}\|f\|_\infty$
than the constant one. An upper bound on such a susceptibility by the variance of the sum of
$\ReTr Q_p$ over the six plaquettes with lower corner $0$ is a consequence
of the averaged condition by Proposition~\ref{rD:suscept}.
\end{remark}

\emph{The free-energy normalisation.} The function
\[
H\mapsto a^{-4}\max_\pm R_M(\beta;\pm H)
\]
is continuous, vanishes at $0$ and is strictly increasing
on $[0,\infty)$. Let
\begin{align*}
h_{M,\beta}&=\sup\bigl\{H>0:\ a^{-4}\max_\pm R_M(\beta;\pm H)\le1\bigr\}\in
(0,\infty],\\
Z^{\mathrm{FE}}_{M,\beta}&=\frac N{4a^4}\min(h_{M,\beta},1).
\end{align*}
On the torus $h_{M,\beta}<\infty$: as $H\to\infty$, $p_M'(\beta+H)\to6N^2$,
because $\mu_{M,\beta+H}$ concentrates on the maximum $S=6NM^4$ of $S$,
attained on the flat configurations, while $p_M'(\beta)<6N^2$; so
$R_M(\beta;H)\ge\int_{H/2}^H(p_M'(\beta+s)-p_M'(\beta))\,ds\to\infty$. The
minimum with $1$ is only needed for the infinite-lattice version of
Sec.~\ref{rU7:infinite}. With this choice,
\begin{equation}
\begin{split}
\Xi_{M,\beta,Z^{\mathrm{FE}}}(T)&=a^{-4}\max_\pm R_M\bigl(\beta;\pm T\min(h_{M,
\beta},1)\bigr)\\
&\le T\qquad(0\le T\le1)
\end{split}
\tag{rU7.2}\label{rU7:xi}
\end{equation}
by $R_M(\beta;sh)\le sR_M(\beta;h)$, $s\in[0,1]$, and the definition of
$h_{M,\beta}$.

\begin{corollary}[unconditional bounds on the torus]\label{rU7:fe}
Let $M_k\ge4$ be even, $a_k=L/M_k\to0$, $\beta_k\ge0$ arbitrary, $Z_k=
Z^{\mathrm{FE}}_{M_k,\beta_k}$ and $\mathcal A_k$ the corresponding smeared
field. Then for every bounded $f\ne0$ on $\mathbb T_L$, every $k$ and every
$t$ with $|t|\,\|f\|_\infty\le1$,
\[
\log\mu_k\bigl(e^{t\mathcal A_k(f)}\bigr)\le|t|\,\|f\|_{1,a_k}\le|t|\,
\|f\|_\infty L^4 ,
\]
and for every $n\ge1$
\begin{align*}
\mu_k\bigl(\mathcal A_k(f)^2\bigr)&\le2e^{L^4}\|f\|_\infty\|f\|_{1,a_k},\\
\mu_k\bigl(|\mathcal A_k(f)|^n\bigr)&\le2\,n!\,e^n\bigl(\|f\|_\infty+\tfrac1n
\|f\|_{1,a_k}\bigr)^n,\\
\mu_k\bigl(|\mathcal A_k(f)|^n\bigr)&\le2\,n!\,e^{L^4}\|f\|_\infty^n,\\
\Bigl|\mu_k\Bigl(\prod_{j=1}^n\mathcal A_k(f_j)\Bigr)\Bigr|&\le2\,n!\,e^{L^4}
\prod_{j=1}^n\|f_j\|_\infty .
\end{align*}
\end{corollary}

\begin{proof}
Theorem~\ref{rU7:torus}(i) and~\eqref{rU7:xi} give the first bound. For
$0<t\le1/\|f\|_\infty$ and $y\in\R$, $|y|^n\le n!\,t^{-n}(e^{ty}+e^{-ty})$
and $t^2y^2\le e^{ty}+e^{-ty}-2$; with the first bound,
$\mu_k(|\mathcal A_k(f)|^n)\le2\,n!\,t^{-n}e^{t\|f\|_{1,a_k}}$ and
$\mu_k(\mathcal A_k(f)^2)\le2t^{-2}(e^{t\|f\|_{1,a_k}}-1)$. The choice
$t=1/\|f\|_\infty$ with $e^u-1\le ue^u$ and $\|f\|_{1,a_k}\le L^4\|f\|_\infty$
gives the second-moment bound and the third bound; the choice
$t=(\|f\|_\infty+\|f\|_{1,a_k}/n)^{-1}$, for which $t\|f\|_{1,a_k}\le n$,
gives the second. The
product bound follows by H\"older's inequality, $|\mu(\prod_j\mathcal A(f_j))|
\le\prod_j\mu(|\mathcal A(f_j)|^n)^{1/n}$.
\end{proof}

\begin{corollary}[the variance normalisation]\label{rU7:var}
With $Z^{\mathrm V}_{M,\beta}=N/(4a^2p_M''(\beta)^{1/2})$ and $\mathcal A=
\mathcal A_{M,\beta,Z^{\mathrm V}_{M,\beta}}$,
\[
\mu_{M,\beta}\bigl(\mathcal A(1)^2\bigr)=L^4,\qquad
\mu_{M,\beta}\bigl(\mathcal A(f)^2\bigr)\le\|f\|_{1,a}\|f\|_\infty
\]
for every even $M\ge4$, every $\beta\ge0$ and every bounded $f$ on
$\mathbb T_L$.
\end{corollary}

\begin{proof}
Theorem~\ref{rU7:torus}(iii) with $16Z^2a^4p_M''(\beta)/N^2=1$.
\end{proof}

The variance normalisation gives second moments only; by
Theorem~\ref{rU7:torus}(iv) it gives exponential moments exactly when the
remainders $a^{-4}R_M(\beta;\pm Ta^2p_M''(\beta)^{-1/2})$ stay bounded, which
holds, for instance, if $p_M''$ varies by a bounded factor on the coupling
window $|\beta'-\beta|\le Ta^2p_M''(\beta)^{-1/2}$, since then $a^{-4}R_M
\le\frac12T^2\sup p_M''/p_M''(\beta)$.
\subsubsection{The infinite lattice}\label{rU7:infinite}

The torus bounds pass to the torus limits of Lemma~\ref{rE:torus} (there
the torus has side $L$; here the side is $M$ and $L$ is the physical length).
The free energy is replaced by a limit of free energies and its derivative at
$\beta$ by a subgradient.

\begin{lemma}[free energy of a torus limit]\label{rU7:pinf}
Let $\beta\ge0$ and let $(M_j)$ be a sequence of even integers along which the
torus measures $\mu_{M_j,\beta}$ converge to a state $\omega=\omega_\beta$ as
in Lemma~\ref{rE:torus}\textup{(i)}. Then $(M_j)$ has a subsequence, again
denoted $(M_j)$, along which $p_{M_j}(\beta')\to p^\infty(\beta')$ locally
uniformly in $\beta'\in\R$ and $p_{M_j}'(\beta)\to D$, where $p^\infty$ is
convex with $|p^\infty(\beta')-p^\infty(\beta'')|\le6N^2|\beta'-\beta''|$,
$D=6N\,\omega(\ReTr Q_p)$, and $D$ is a subgradient of $p^\infty$ at
$\beta$. Consequently
\[
R^\infty(h)=p^\infty(\beta+h)-p^\infty(\beta)-hD\ \ge0\qquad(h\in\R)
\]
is convex, $R^\infty(0)=0$, $R^\infty(sh)\le sR^\infty(h)$ for $s\in[0,1]$,
and $R^\infty$ is nondecreasing on $[0,\infty)$ and nonincreasing on
$(-\infty,0]$. Moreover $c_{M_j}=\mu_{M_j,\beta}(\Phi^\square_0)\to
c^\infty=\omega(\Phi^\square_0)=24(N-D/(6N))$, and for every finitely
supported $\lambda:\Z^4\to\R$,
\[
\omega\Bigl(\exp\sum_{x\in\Z^4}\lambda_x\bigl(\Phi^\square_x-c^\infty\bigr)
\Bigr)\le\exp\Bigl(\sum_{x\in\Z^4}R^\infty\bigl(-4\lambda_x/N\bigr)\Bigr).
\]
\end{lemma}

\begin{proof}
Each $p_M$ is convex with $p_M(0)=0$ and $|p_M'|\le6N^2$, so the family
$(p_{M})$ is uniformly bounded and equi-Lipschitz on compact sets; by
Arzel\`a--Ascoli and a diagonal argument a subsequence converges locally
uniformly to a function $p^\infty$, which is convex and $6N^2$-Lipschitz as a
locally uniform limit of such functions. Passing to a further subsequence,
$p_{M_j}'(\beta)=6N\mu_{M_j,\beta}(\ReTr Q_p)$ converges, and its limit is
$6N\omega(\ReTr Q_p)=D$ because $\ReTr Q_p$ is a continuous function of the four
links of $p$; here $p_{M}'(\beta)=6N\mu_{M,\beta}(\ReTr Q_p)$ holds for each of
the six plaquettes $p$ with lower corner $0$, and so does the limit, by
the $B_4$ invariance of $\omega$ (Lemma~\ref{rE:torus}(ii)). For each $j$ and $h$, $p_{M_j}(\beta+h)\ge p_{M_j}(\beta)+hp_{M_j}'(
\beta)$ by convexity; letting $j\to\infty$ gives $p^\infty(\beta+h)\ge
p^\infty(\beta)+hD$, that is, $D$ is a subgradient and $R^\infty\ge0$. The
remaining properties of $R^\infty$ follow from convexity and $R^\infty(0)=0$
exactly as for $R_M$. The convergence $c_{M_j}\to c^\infty$ holds because
$\Phi^\square_0$ is a continuous function of finitely many links, and
$c_{M_j}=24(N-p_{M_j}'(\beta)/(6N))$.

For the last bound fix $\lambda$ with support in a box of side $2m$ centred at
the origin. For $M_j>2m+4$ the box embeds in $\Lambda_{M_j}$ and
Lemma~\ref{rU7:chess} gives $\mu_{M_j,\beta}(\exp\sum_x\lambda_x(\Phi^\square
_x-c_{M_j}))\le\exp(\sum_xR_{M_j}(\beta;-4\lambda_x/N))$. The function
$\exp\sum_x\lambda_x(\Phi^\square_x-c)$ is continuous in finitely many links
and jointly continuous in $c$, and $\exp\sum_x\lambda_x(\Phi^\square_x-c_{M_j})
-\exp\sum_x\lambda_x(\Phi^\square_x-c^\infty)$ tends to $0$ uniformly on the
configuration space, so the left side converges to the left side of the claim.
On the right, $R_{M_j}(\beta;h)\to R^\infty(h)$ for each $h$ by the locally
uniform convergence of $p_{M_j}$ and the convergence of $p_{M_j}'(\beta)$.
\end{proof}

\emph{The sequence of states.} Let $\beta_k\ge0$ be arbitrary and
$a_k\to0$, $a_k\le1$. For each $k$ fix a torus limit $\omega_k$ of
$\mu_{M,\beta_k}$ along a sequence of even $M$ for which also the conclusions
of Lemma~\ref{rU7:pinf} hold, with $p_k^\infty$, $D_k$, $R_k^\infty$ and
$c_k=\omega_k(\Phi^\square_0)$. For a bounded function $f$ on $\R^4$ with
compact support, $Z>0$ and $m_x=x+\frac12\mathbf1$, put
\begin{align*}
A_{k,Z}(f)&=Z\sum_{x\in\Z^4}a_k^4f(a_km_x)\bigl(\Phi^\square_x-c_k\bigr),\\
\|f\|_{1,k}&=\sum_{x\in\Z^4}a_k^4|f(a_km_x)| ,\qquad
\mathfrak m_k(f)=\frac{\|f\|_{1,k}}{\|f\|_\infty}\quad(f\ne0),
\end{align*}
a finite sum. If $f$ vanishes outside a closed cube $Q$ of side $\ell$, then
at most $(\ell/a_k+1)^4$ points of the lattice $a_k(\Z^4+\frac12\mathbf1)$
lie in $Q$, so
\[
\mathfrak m_k(f)\le(\ell+a_k)^4\le(\ell+1)^4 .
\tag{rU7.3}\label{rU7:count}
\]
For $f$ Riemann integrable, $\|f\|_{1,k}\to\|f\|_{L^1}$ as $k\to\infty$,
since $\|f\|_{1,k}$ is a Riemann sum of $|f|$ over the cubes of side $a_k$
centred at the points $a_km_x$.

\begin{theorem}[chessboard bound and free-energy normalisation on $\Z^4$]
\label{rU7:zfour}
In the setting above let
\begin{align*}
\Xi_k^\infty(T)&=a_k^{-4}\max_\pm R_k^\infty\bigl(\pm4TZ_ka_k^4/N\bigr),\\
h_k^\infty&=\sup\bigl\{H>0:\ a_k^{-4}\max_\pm R_k^\infty(\pm H)\le1\bigr\}\in
(0,\infty],
\end{align*}
and $Z_k^{\mathrm{FE}}=\frac N{4a_k^4}\min(h_k^\infty,1)$.

\textup{(i)} For every $Z_k>0$, every bounded $f\ne0$ with compact support and
every $t\in\R$,
$\log\omega_k(e^{tA_{k,Z_k}(f)})\le\mathfrak m_k(f)\,\Xi_k^\infty(|t|\,\|f\|
_\infty)$.

\textup{(ii)} With $Z_k=Z_k^{\mathrm{FE}}$, $\Xi_k^\infty(T)\le T$ for $0\le T
\le1$; hence, writing $A_k=A_{k,Z_k^{\mathrm{FE}}}$, for $|t|\,\|f\|_\infty
\le1$
\[
\log\omega_k\bigl(e^{tA_k(f)}\bigr)\le|t|\,\|f\|_{1,k},
\]
and for every $n\ge1$ and every $f,f_1,\dots,f_n$ vanishing outside a cube of
side $\ell$,
\begin{align*}
\omega_k\bigl(|A_k(f)|^n\bigr)&\le2\,n!\,e^{(\ell+1)^4}\|f\|_\infty^n,\\
\omega_k\bigl(A_k(f)^2\bigr)&\le2e^{(\ell+1)^4}\,\|f\|_\infty\|f\|_{1,k},
\end{align*}
\[
\Bigl|\omega_k\Bigl(\prod_{j=1}^nA_k(f_j)\Bigr)\Bigr|\le2\,n!\,e^{(\ell+1)^4}
\prod_{j=1}^n\|f_j\|_\infty .
\]
These bounds hold for every $k$, with no condition on $\beta_k$.
\end{theorem}

\begin{proof}
(i) is the last bound of Lemma~\ref{rU7:pinf} with $\lambda_x=tZ_ka_k^4f(a_k
m_x)$, followed by the convexity argument in the proof of
Theorem~\ref{rU7:torus}(i), which used only $R(0)=0$ and convexity.

(ii) If $h_k^\infty<\infty$ then $a_k^{-4}\max_\pm R_k^\infty(\pm h_k^\infty)
\le1$ by continuity of $R_k^\infty$, and $4TZ_k^{\mathrm{FE}}a_k^4/N=T\min(h_k
^\infty,1)\le Th_k^\infty$, so $\Xi_k^\infty(T)\le T\,a_k^{-4}\max_\pm
R_k^\infty(\pm h_k^\infty)\le T$ for $T\le1$ by $R^\infty(sh)\le sR^\infty
(h)$. If $h_k^\infty=\infty$ then $a_k^{-4}\max_\pm R_k^\infty(\pm H)\le1$
for all $H$ and the same argument with $H=1$ applies. The exponential bound
follows from (i). For the moments, $|y|^n/n!\le T^{-n}(e^{Ty}+e^{-Ty})$ with
$T=1/\|f\|_\infty$ and $\mathfrak m_k(f)\le(\ell+1)^4$ give the first bound.
For the second, $t^2y^2\le e^{ty}+e^{-ty}-2$ with $t=1/\|f\|_\infty$ gives
$\omega_k(A_k(f)^2)\le\|f\|_\infty^2\cdot2(e^{\|f\|_{1,k}/\|f\|_\infty}-1)
\le2\|f\|_\infty\|f\|_{1,k}e^{\mathfrak m_k(f)}$, using $e^u-1\le ue^u$. The
product bound is H\"older's inequality.
\end{proof}

\begin{remark}\label{rU7:zfourremark}
(a) Theorem~\ref{rU7:zfour}(ii) is a compactness statement for the family of
all smeared, centred plaquette fields along \emph{any} sequence of
couplings: the normalisation $Z_k^{\mathrm{FE}}$ is read off from the
infinite-volume free energy alone, and with it all moments are bounded
uniformly in $k$, including at coinciding points and in the ultraviolet
limit, by constants depending only on the size of the support and on
$\|f\|_\infty$. It assumes nothing about decay, about the tuning or about the
averaged condition. Its price is paid in Sec.~\ref{rU7:af}: when the
short-distance law (AF$^\square$) holds for constants $Z_k'$,
$Z_k^{\mathrm{FE}}/Z_k'\to0$, and the limits obtained with
$Z_k^{\mathrm{FE}}$ have vanishing separated two-point functions
(Theorem~\ref{rU7:afthm}).

(b) The second-moment bound is of the form $\|f\|_\infty\|f\|_{1,k}$, not
$\|f\|_\infty^2$; this two-norm form is what makes the limit field continuous
under translations and reflections of step functions
(Theorem~\ref{rU7:limit}).

(c) By Theorem~\ref{rU7:torus}(ii), for the torus states the choice
$Z^{\mathrm{FE}}$ is sharp up to the factor $\min(h,1)$: with any $Z_k\ge
TZ_k^{\mathrm{FE}}$, $T>1$, and $h_k\le1$, the exponential moment of
$\pm\mathcal A_k(1)$ at $t=1$ is at least $e^{L^4}$ in one of the signs, and
with $Z_k/Z_k^{\mathrm{FE}}\to\infty$ it diverges.
\end{remark}

\emph{Test functions.} Let $\mathcal R$ be the space of bounded, compactly
supported, Riemann integrable functions $f:\R^4\to\R$ (bounded, compactly
supported, continuous outside a Lebesgue null set); it is an algebra
containing $C_c(\R^4)$ and the indicators of boxes, and it is stable under
translations $\tau_vf=f(\cdot-v)$, under the hypercubic group $B_4$ acting by
$f\mapsto f\circ R^{-1}$, and under the reflections $\theta_{\mu,c}$ in the
planes $\{x_\mu=c\}$. Let $\mathcal G\subset\mathcal R$ be the countable
$\mathbb Q$-vector space spanned by the indicators of boxes $\prod_\mu[a_\mu,b_\mu)$
with rational corners.

\begin{lemma}[approximation]\label{rU7:approx}
Let $f\in\mathcal R$ vanish outside a cube $Q$ of side $\ell$ and $\varepsilon
\in(0,1]$. There is $g\in\mathcal G$ vanishing outside a cube of side $\ell+1$
containing $Q$ with $\|g\|_\infty\le\|f\|_\infty+1$ and $\|f-g\|_{L^1}<
\varepsilon$. For such $f,g$ and all large $k$,
\[
\omega_k\bigl(A_k(f-g)^2\bigr)\le4\varepsilon\,(2\|f\|_\infty+1)\,
e^{(\ell+2)^4}.
\]
Moreover, for $g\in\mathcal G$ and $v,w\in\R^4$ with $|v-w|_\infty\le2a_k$,
\[
\omega_k\bigl(A_k(\tau_vg-\tau_wg)^2\bigr)\le C_g\,a_k ,
\]
with $C_g$ depending only on $g$.
\end{lemma}

\begin{proof}
By the Darboux criterion there is a partition of a cube $Q'\supset Q$ of side
$\ell+1$ with rational corners into boxes $B_i$ such that $\sum_i(\sup_{B_i}f
-\inf_{B_i}f)|B_i|<\varepsilon/2$; put $g=\sum_iq_i\mathbf 1_{B_i}$ with
rational $q_i\in[\inf_{B_i}f,\inf_{B_i}f+\varepsilon/(2|Q'|)]$. Then
$\|g\|_\infty\le\|f\|_\infty+1$ (as $|Q'|\ge1$) and $\|f-g\|_{L^1}\le\sum_i
(\sup_{B_i}f-\inf_{B_i}f)|B_i|+\varepsilon/2<\varepsilon$. (Such a partition
exists: $f$ is Riemann integrable on $Q'$, so the upper and lower Darboux sums
of every sufficiently fine grid partition differ by less than $\varepsilon/2$,
and the grid points can be taken rational.) Since $f-g\in\mathcal R$,
$\|f-g\|_{1,k}\to\|f-g\|_{L^1}$, so $\|f-g\|_{1,k}<2\varepsilon$ for large
$k$, and the second bound of Theorem~\ref{rU7:zfour}(ii) with
$\|f-g\|_\infty\le2\|f\|_\infty+1$ and the cube of side $\ell+1$ gives the
claim. For the last bound write $g=\sum_iq_i\mathbf 1_{B_i}$ with boxes $B_i$
contained in a cube of side $\ell_g$, let $n_g$ be the total number of
$3$-dimensional faces of the $B_i$, and put $r=|v-w|_\infty\le2a_k\le2$.
The function $h=\tau_vg-\tau_wg$ is bounded by $2\|g\|_\infty$ and vanishes
outside a cube of side $\ell_g+2$. If $h(y)\ne0$, then for some $i$ exactly
one of $y-v$, $y-w$ lies in $B_i$; the segment between them has
$\ell^\infty$-length $r$, so $y-v$ lies within $\ell^\infty$-distance $2r$
of some face $F$ of $B_i$, that is, $y\in v+F+[-2r,2r]^4$. This set is a
box with one side at most $4r\le8a_k$ and three sides at most $\ell_g+8a_k$,
so it contains at most $9(\ell_g/a_k+9)^3$ points of the lattice
$a_k(\Z^4+\frac12\mathbf 1)$. Hence, with $a_k\le1$,
$\|h\|_{1,k}\le2\|g\|_\infty\cdot9n_ga_k^4(\ell_g/a_k+9)^3\le18n_g\|g\|_\infty
(\ell_g+9)^3a_k$, and the second-moment bound of
Theorem~\ref{rU7:zfour}(ii), with a cube of side $\ell_g+2$, gives the claim
with $C_g=72n_g\|g\|_\infty^2(\ell_g+9)^3e^{(\ell_g+3)^4}$.
\end{proof}
\subsubsection{Subsequential continuum limits}\label{rU7:limits}

We treat the torus of side $L$ and the infinite lattice together. In the
\emph{torus setting} $X=\mathbb T_L$, $\nu_k=\mu_{M_k,\beta_k}$ with $M_k\ge4$
even, $a_k=L/M_k\to0$, $A_k=\mathcal A_{M_k,\beta_k,Z^{\mathrm{FE}}_{M_k,
\beta_k}}$, $\|f\|_{1,k}=\|f\|_{1,a_k}$, $\mathcal R(X)$ is the space of
bounded functions on $\mathbb T_L$ continuous outside a null set, and for a
direction $\mu$ and $c\in\R$ the \emph{positive region} is $P_{\mu,c}=\{c<
x_\mu<c+\frac L2\}$ (modulo $L$). In the \emph{infinite-lattice setting}
$X=\R^4$, $\nu_k=\omega_k$, $A_k=A_{k,Z^{\mathrm{FE}}_k}$ and $\|f\|_{1,k}$ are
those of Theorem~\ref{rU7:zfour}, $\mathcal R(X)=\mathcal R$, and $P_{\mu,c}=
\{x_\mu>c\}$. In both, $\beta_k\ge0$ is arbitrary, $\theta_{\mu,c}$ is the
reflection $x_\mu\mapsto2c-x_\mu$, $\theta f=f\circ\theta$, $\tau_vf=f(\cdot-v)$,
$B_4$ acts on $X$ fixing $0$, and $\Gamma(X)$ is the group generated by the
translations of $X$ and $B_4$. The bounds of Corollary~\ref{rU7:fe}, and in the
infinite-lattice setting those of Theorem~\ref{rU7:zfour}(ii) together with
the bound $\nu_k(|A_k(f)|^n)\le2\,n!\,e^n(\|f\|_\infty+\|f\|_{1,k}/n)^n$, which
follows from Theorem~\ref{rU7:zfour}(ii) exactly as in the proof of
Corollary~\ref{rU7:fe}, hold in both settings. On $\mathbb T_L$ we let
$\mathcal G$ be the $\mathbb Q$-span of the indicators of the images in
$\mathbb T_L$ of boxes $\prod_\mu[a_\mu,b_\mu)$ with $0\le b_\mu-a_\mu\le L$ and
corners in the countable dense set $\mathbb Q+L\mathbb Q$, which contains
$0$ and $L$; then Lemma~\ref{rU7:approx} holds on $\mathbb T_L$ with the same
proof, the cube $Q'$ being a fundamental domain $[0,L)^4$ when $\ell+1>L$,
and with $e^{(\ell+2)^4}$, $e^{(\ell_g+3)^4}$ replaced by $e^{L^4}$
(Corollary~\ref{rU7:fe}).

\begin{lemma}[exact lattice symmetries]\label{rU7:latsym}
\textup{(a)} For $\gamma=\tau_v$ with $v\in a_k\Z^4$, or $\gamma\in B_4$, the
law under $\nu_k$ of $(A_k(f_j\circ\gamma^{-1}))_{j\le n}$ equals that of
$(A_k(f_j))_{j\le n}$, for all $f_j\in\mathcal R(X)$.

\textup{(b)} Let $c\in a_k\Z$, let $f_1,\dots,f_n\in\mathcal R(X)$ vanish outside
a compact subset of $P_{\mu,c}$, and let $F:\R^n\to\C$ be bounded and
continuous. Then, with $\theta=\theta_{\mu,c}$,
\[
\nu_k\Bigl(\overline{F\bigl(A_k(\theta f_1),\dots,A_k(\theta f_n)\bigr)}\,
F\bigl(A_k(f_1),\dots,A_k(f_n)\bigr)\Bigr)\ge0 .
\]
\end{lemma}

\begin{proof}
(a) The translation of the lattice by $v/a_k\in\Z^4$ maps $C_x$ onto
$C_{x+v/a_k}$ and $m_x$ to $m_x+v/a_k$. An element $g\in B_4$ maps
$\Z^4$ to itself and $\frac12\mathbf1$ to a vector with entries $\pm\frac12$,
so it maps $m_x$ to $m_y$ with $y=gx+g(\frac12\mathbf1)-\frac12\mathbf1\in\Z^4$,
and the cube $C_x$ onto $C_y$. In both cases, with the action on link
variables of the proof of Lemma~\ref{rE:torus}(ii), $\Phi^\square_x$ is carried
to $\Phi^\square$ of the image cube, because $\ReTr Q^{-1}=\ReTr Q$. Hence
$A_k(f\circ\gamma^{-1})$ is the image of $A_k(f)$ under the corresponding
transformation of configurations, which preserves $\nu_k$ (on the torus
directly, on $\Z^4$ by Lemma~\ref{rE:torus}(ii)); the constant $c_k$ does not
depend on $x$.

(b) Let $\vartheta$ be the lattice reflection in $\{x_\mu=c/a_k\}$. If
$f_j(a_km_x)\ne0$ then $a_km_x\in P_{\mu,c}$; since $m_{x,\mu}-c/a_k$ is a
half-integer, it is at least $\frac12$, so $x_\mu\ge c/a_k$, and on the torus
also $x_\mu+1\le c/a_k+M_k/2$. So $A_k(f_j)$ is a function of the links of
cubes in the closed half on the positive side of $\vartheta$, and
$\vartheta A_k(f_j)=A_k(\theta f_j)$, since $\vartheta$ maps $C_x$ onto the
cube with centre $\vartheta m_x$ and $\theta(a_km_x)=a_k\vartheta(m_x)$. On
the torus the claim is the reflection positivity proved in
Lemma~\ref{rU7:chess} for the continuous function $F(A_k(f_1),\dots)$ of the
links of the closed half. On $\Z^4$ the variables $A_k(f_j)$ depend on
finitely many links; for each large $M$ in the sequence defining $\omega_k$
the same inequality holds for $\mu_{M,\beta_k}$ with the reflection of
$\Lambda_M$ in $\{x_\mu=c/a_k\}$, and it passes to the limit.
\end{proof}

\begin{theorem}[subsequential limits of the smeared plaquette field]
\label{rU7:limit}
In either setting, every sequence of $k$ has a subsequence along which the
following holds: there are a probability space $(\Omega,\mathcal F,\mathrm P)$
and real random variables $A(f)$, $f\in\mathcal R(X)$, such that

\textup{(i)} for all $f_1,\dots,f_n\in\mathcal R(X)$, $(A_k(f_j))_j\to(A(f_j))
_j$ in law, $\nu_k(\prod_jA_k(f_j))\to\mathrm E\prod_jA(f_j)$, and
$\nu_k(e^{tA_k(f)})\to\mathrm Ee^{tA(f)}$ for $|t|\,\|f\|_\infty<1$;

\textup{(ii)} $f\mapsto A(f)$ is linear almost surely, $\mathrm EA(f)=0$, and
for $|t|\,\|f\|_\infty\le1$ and $n\ge1$
\begin{align*}
\log\mathrm Ee^{tA(f)}&\le|t|\,\|f\|_{L^1},\\
\mathrm E|A(f)|^n&\le2\,n!\,e^n\bigl(\|f\|_\infty+\tfrac1n\|f\|_{L^1}\bigr)^n,\\
\mathrm EA(f)^2&\le2\|f\|_\infty\|f\|_{L^1}e^{\|f\|_{L^1}/\|f\|_\infty},\\
\Bigl|\mathrm E\prod_{j=1}^nA(f_j)\Bigr|&\le2\,n!\,e^n\prod_{j=1}^n\bigl(
\|f_j\|_\infty+\|f_j\|_{L^1}\bigr);
\end{align*}

\textup{(iii)} for every $\gamma\in\Gamma(X)$ the laws of
$(A(f_j\circ\gamma^{-1}))_j$ and $(A(f_j))_j$ coincide, and $v\mapsto
A(\tau_vf)$ is continuous into $L^2(\mathrm P)$;

\textup{(iv)} for every direction $\mu$, every $c\in\R$, all $f_1,\dots,f_n\in
\mathcal R(X)$ vanishing outside a compact subset of $P_{\mu,c}$ and every
bounded continuous $F$, with $\theta=\theta_{\mu,c}$,
\[
\mathrm E\Bigl[\overline{F\bigl(A(\theta f_1),\dots,A(\theta f_n)\bigr)}\,
F\bigl(A(f_1),\dots,A(f_n)\bigr)\Bigr]\ge0 ;
\]
in particular $\sum_{i,j}\bar\alpha_i\alpha_j\,\mathrm Ee^{i(A(f_j)-A(\theta
f_i))}\ge0$ for $\alpha_i\in\C$, and the same holds with $F$ a polynomial.

\textup{(v)} In the torus setting, with $\eta_k=\min(h_{M_k,\beta_k},1)$,
for $|t|<1$
\begin{align*}
\log\mathrm Ee^{\pm tA(1)}&=\lim_kL^4a_k^{-4}R_{M_k}(\beta_k;\mp t\eta_k),\\
\mathrm EA(1)^2&=L^4\lim_k\eta_k^2\,p_{M_k}''(\beta_k)/a_k^4,
\end{align*}
the limits existing along the subsequence; so $A\ne0$ if and only if this
last limit is positive.
\end{theorem}

\begin{proof}
Let $f\in\mathcal R(X)$. The bounds quoted before Lemma~\ref{rU7:latsym}
give
\[
\sup_k\nu_k(|A_k(f)|^n)<\infty\quad(n\ge1),\qquad
\sup_k\nu_k\bigl(e^{|A_k(f)|/\|f\|_\infty}\bigr)<\infty
\]
(in the infinite-lattice
setting with constants depending on a cube containing the support); hence,
for each finite family, the products $\prod_jA_k(f_j)$ and the variables
$e^{tA_k(f)}$, $|t|\,\|f\|_\infty<1$, are uniformly integrable, by H\"older's
inequality and because $e^{|t||y|}$ is a power $<1$ of $e^{|y|/\|f\|
_\infty}$.

\emph{Step 1 (a countable family).} Let $\mathcal G$ be the countable space of
rational step functions of Lemma~\ref{rU7:approx}. For each finite family in
$\mathcal G$ the laws of $(A_k(g_j))_j$ on $\R^n$ are tight, by the second
moments. A diagonal argument over the countably many finite families gives a
subsequence along which all these laws converge; the limits are consistent
under projections, and the Kolmogorov extension theorem gives random
variables $(A(g))_{g\in\mathcal G}$ with these finite-dimensional laws. For
rational $q,q'$, $A_k(qg+q'g')=qA_k(g)+q'A_k(g')$ identically, so the law of
$(A_k(qg+q'g'),A_k(g),A_k(g'))$ is carried by a closed hyperplane, and so is
its limit: $A(qg+q'g')=qA(g)+q'A(g')$ almost surely.

\emph{Step 2 (extension).} Let $f\in\mathcal R(X)$ and let $g_m\in\mathcal G$ be
as in Lemma~\ref{rU7:approx} with $\|f-g_m\|_{L^1}\to0$. By the second-moment
bound, with a constant $C_f$ depending only on $\|f\|_\infty$ and on a cube
containing the supports,
\begin{align*}
\limsup_k\nu_k\bigl(A_k(f-g_m)^2\bigr)&\le C_f\|f-g_m\|_{L^1},\\
\mathrm E\bigl(A(g_m)-A(g_{m'})\bigr)^2&\le C_f\|g_m-g_{m'}\|_{L^1},
\end{align*}
the second because $\mathrm E(A(g)-A(g'))^2=\lim_k\nu_k(A_k(g-g')^2)$ by
uniform integrability and $\|g-g'\|_{1,k}\to\|g-g'\|_{L^1}$. So $A(g_m)$ is
Cauchy in $L^2(\mathrm P)$; let $A(f)$ be its limit. The same two bounds
show that the limit is the same for every sequence $g_m\in\mathcal G$ with
$\sup_m\|g_m\|_\infty<\infty$, supports in a fixed cube and $\|f-g_m\|_{L^1}
\to0$, and that it agrees with Step~1 for $f\in\mathcal G$. For $f_1,\dots,f_n$ and $t\in\R^n$, using $|e^{iy}-e^{iy'}|\le
|y-y'|$,
\begin{align*}
&\Bigl|\nu_k\bigl(e^{i\sum_jt_jA_k(f_j)}\bigr)-\mathrm Ee^{i\sum_jt_jA(f_j)}
\Bigr|\\
&\quad\le\Bigl|\nu_k\bigl(e^{i\sum_jt_jA_k(g_{j,m})}\bigr)-\mathrm Ee^{i\sum
_jt_jA(g_{j,m})}\Bigr|\\
&\qquad+\sum_j|t_j|\bigl(\delta_{j,m,k}+\delta_{j,m}\bigr),
\end{align*}
with $\delta_{j,m,k}=\nu_k(A_k(f_j-g_{j,m})^2)^{1/2}$ and $\delta_{j,m}=
\mathrm E((A(f_j)-A(g_{j,m}))^2)^{1/2}$; letting $k\to\infty$ and then
$m\to\infty$ gives convergence in law. Convergence of the moments and of the
exponential moments follows by uniform integrability. For linearity let
$f,f'\in\mathcal R(X)$, $\lambda,\lambda'\in\R$, let $g_m,g_m'$ be sequences
as above for $f,f'$, and let $q_m,q'_m\in\mathbb Q$ tend to $\lambda,
\lambda'$; then $q_mg_m+q'_mg'_m\in\mathcal G$ is such a sequence for
$\lambda f+\lambda'f'$, and $A(q_mg_m+q'_mg'_m)=q_mA(g_m)+q'_mA(g'_m)$ almost
surely by Step~1, so $A(\lambda f+\lambda'f')=\lambda A(f)+\lambda'A(f')$
almost surely. (Linearity is meant in this sense: almost surely for each
finite linear relation.) Finally $\mathrm EA(f)=\lim_k\nu_k(A_k(f))=0$.

\emph{Step 3 (bounds).} For continuous $\varphi\ge0$ and convergence in law,
$\mathrm E\varphi(A(f))\le\liminf_k\nu_k(\varphi(A_k(f)))$; applied to
$e^{ty}$, $|y|^n$, $y^2$ with the bounds of Corollary~\ref{rU7:fe} or
Theorem~\ref{rU7:zfour}(ii) and $\|f\|_{1,k}\to\|f\|_{L^1}$, this gives the first
three bounds of (ii); the last is H\"older's inequality.

\emph{Step 4 (symmetries).} For $\gamma\in B_4$, Lemma~\ref{rU7:latsym}(a)
holds for every $k$ and passes to the limit in law. For $v\in X$ let $v_k\in
a_k\Z^4$ with $|v-v_k|_\infty\le a_k$. For $g\in\mathcal G$,
$\nu_k(A_k(\tau_vg-\tau_{v_k}g)^2)\le C_ga_k$ by Lemma~\ref{rU7:approx}, and the
law of $(A_k(\tau_{v_k}g_j))_j$ is that of $(A_k(g_j))_j$; so the limit laws of
$(A(\tau_vg_j))_j$ and $(A(g_j))_j$ coincide. Since $\tau_vg_m$ approximates
$\tau_vf$ as $g_m$ approximates $f$, this extends to $\mathcal R(X)$. Continuity: for $f\in\mathcal R(X)$ and $h=\tau_vf-f$, $|v|\le1$, the third
bound of (ii) (proved in Step~3, which does not use Step~4) and
$\|h\|_{L^1}\le\|h\|_\infty|\{h\ne0\}|$ give $\mathrm E(A(\tau_vf)-A(f))^2=
\mathrm EA(h)^2\le4\|f\|_\infty e^{K}\|\tau_vf-f\|_{L^1}\to0$ as $v\to0$, with
$K$ the volume of the $1$-neighbourhood of the support of $f$ (or
$K=L^4$ on the torus).

\emph{Step 5 (positivity).} Since $0\in a_k\Z$ for every $k$,
Lemma~\ref{rU7:latsym}(b) with $c=0$ holds for every $k$; the function
$(y,y')\mapsto\overline{F(y')}F(y)$ is bounded and continuous on $\R^{2n}$, and
the joint law of $(A_k(f_j),A_k(\theta_{\mu,0}f_j))_j$ converges, so the
inequality passes to the limit. For general $c$ put $g_j=f_j(\cdot+ce_\mu)$,
which vanish outside a compact subset of $P_{\mu,0}$; since $\theta_{\mu,c}x-
ce_\mu=\theta_{\mu,0}(x-ce_\mu)$, the family $(f_j,\theta_{\mu,c}f_j)_j$ is the
translate by $ce_\mu$ of $(g_j,\theta_{\mu,0}g_j)_j$, and (iii) reduces the
claim to $c=0$. For a polynomial $F$, truncate: $F_R=F\chi(|y|/R)$ with a
continuous cut-off $\chi$, and let $R\to\infty$ using the moment bounds.

\emph{Step 6 (the constant smearing).} By Theorem~\ref{rU7:torus}(ii) and
$4Z_k^{\mathrm{FE}}a_k^4/N=\eta_k$,
\[
\log\mu_k\bigl(e^{\pm tA_k(1)}\bigr)=L^4a_k^{-4}R_{M_k}(\beta_k;\mp t\eta_k),
\]
and by (i)
the left side converges for $|t|<1$; $\mu_k(A_k(1)^2)=16(Z^{\mathrm{FE}}_k)^2
a_k^4L^4p''_{M_k}(\beta_k)/N^2$ by the proof of Theorem~\ref{rU7:torus}(iii),
and it converges to $\mathrm EA(1)^2$ by uniform integrability. If
$\mathrm EA(1)^2>0$ then $A\ne0$; if $\mathrm EA(1)^2=0$, then $A(1)=0$
and $\mathrm EA(f)^2\le$ the limit of
$\|f\|_{1,a_k}\|f\|_\infty L^{-4}\mu_k(A_k(1)^2)$ by
Theorem~\ref{rU7:torus}(iii), so $A=0$.
\end{proof}

\begin{remark}[what the theorem does and does not give]\label{rU7:limitremark}
Theorem~\ref{rU7:limit} is the compactness half of (E) for the smeared
plaquette field, with no hypothesis on the couplings: a random field on
$\mathbb T_L$ or $\R^4$, invariant under translations and $B_4$, reflection
positive for all coordinate reflections, with moments of factorial growth
bounded by the sup and $L^1$ norms of the test functions; on $\R^4$ these
bounds grow polynomially in the volume of the support. It does not decide
whether $A\ne0$ (on the torus this is the free-energy condition (v)), whether
the limit depends on the subsequence, whether it is $\mathrm O(4)$ invariant,
or whether it is Gaussian. The next subsection shows that, whenever the
short-distance law holds, the normalisation $Z^{\mathrm{FE}}$ is necessarily
too small to see correlations at separated points.
\end{remark}
\subsubsection{The free-energy normalisation against the short-distance law}
\label{rU7:af}

Let $G^\square_k(z)=\nu_k(\Phi^\square_0\Phi^\square_z)-c_k^2$, $z\in\Z^4$, in
either setting of Sec.~\ref{rU7:limits} (on the torus $z\in\Lambda_{M_k}$ and
$|z|$ is the torus distance), so that $\nu_k((\Phi^\square_x-c_k)(\Phi^\square
_y-c_k))=G^\square_k(y-x)$ by translation invariance. With $w_\gamma(\rho)=
\rho^{-8}(1+\log^+\frac1\rho)^{-\gamma}$ as in Sec.~\ref{ssec:limitE}, the
short-distance law for the cube field reads:
\begin{enumerate}[label=\textup{(AF$^\square$)},leftmargin=*]
\item there are $\gamma\ge0$, $0<c_1\le C_1$, $r_0\in(0,\frac14]$ and $k_0$
such that for $k\ge k_0$ and $|z|_\infty\ge2$,
$Z_k'^2|G^\square_k(z)|\le C_1w_\gamma(a_k|z|)$, and
$Z_k'^2G^\square_k(z)\ge c_1w_\gamma(a_k|z|)$ if moreover $a_k|z|\le r_0$.
\end{enumerate}
We call the two inequalities the upper and the lower half. For $|z|_\infty\ge
2$ the cubes $C_0$ and $C_z$ have no common link, which is why $2$ replaces
the $2r+1$ of (AF). The law is the same statement as (AF) for the
gauge-invariant field $\Phi^\square$, centred at the points $m_x$.

\begin{lemma}[two-norm second moments exclude the short-distance law]
\label{rU7:incompat}
Let $Z_k,Z_k'>0$ and suppose that, for some $C<\infty$, all large $k$ and all
$f\in\mathcal R(X)$ vanishing outside a fixed cube $Q_*\subset X$ of side
$\ell_*\in(0,1]$,
\[
\nu_k\bigl(A_{k,Z_k}(f)^2\bigr)\le C\,\|f\|_\infty\|f\|_{1,k},
\]
where $A_{k,Z}$ is the field with normalisation $Z$. If the lower half of
\textup{(AF$^\square$)} holds for the constants $Z_k'$, then $Z_k/Z_k'\to0$.
\end{lemma}

\begin{proof}
Suppose $Z_k/Z'_k\ge\vartheta>0$ along a subsequence. Let $0<s\le r_0/4$ (on $\mathbb T_L$ also $3s\le L/2$, so that the torus
distance of two points of $Q_*$ used below is their Euclidean distance) and
let $B=b+[0,s)^4$ and $B'=2se_0+B$ with $b$ chosen so that $B\cup B'\subset
Q_*$ (possible once $3s\le\ell_*$), $f=\mathbf 1_B$, $g=\mathbf1_{B'}$. If
$a_km_x\in B$ and $a_km_y\in B'$ then $s<|a_km_y-a_km_x|<\sqrt{12}\,s<4s\le
r_0$ and $|y-x|_\infty>s/a_k-1\ge2$ for large $k$. On $[s,4s]$,
$w_\gamma(\rho)\ge(4s)^{-8}(1+\log\frac1s)^{-\gamma}$. So the lower half of
(AF$^\square$) gives, for large $k$,
\begin{align*}
&\nu_k\bigl(A_{k,Z_k}(f)A_{k,Z_k}(g)\bigr)\\
&\quad=Z_k^2\sum_{x,y}a_k^8f(a_km_x)g(a_km_y)G^\square_k(y-x)\\
&\quad\ge\vartheta^2c_1(4s)^{-8}\Bigl(1+\log\frac1s\Bigr)^{-\gamma}
\|f\|_{1,k}\|g\|_{1,k},
\end{align*}
while the hypothesis and the Cauchy--Schwarz inequality bound the left side
by
\[
C(\|f\|_{1,k}\|g\|_{1,k})^{1/2}.
\]
Since $\|f\|_{1,k},\|g\|_{1,k}\to s^4$,
\[
\vartheta^2c_14^{-8}\Bigl(1+\log\frac1s\Bigr)^{-\gamma}\le Cs^4
\]
for every small $s>0$, which fails as $s\to0$.
\end{proof}

\begin{theorem}[what the free energy can see]\label{rU7:afthm}
In either setting of Sec.~\ref{rU7:limits}, let $\beta_k\ge0$ be arbitrary.

\textup{(i)} $Z^{\mathrm{FE}}_k\le N/(4a_k^4)$ for every $k$.

\textup{(ii)} If the lower half of \textup{(AF$^\square$)} holds for constants
$Z_k'$, then $Z_k'\ge c\,a_k^{-4}(1+\log\frac1{a_k})^{-\gamma/2}$ for large $k$,
with $c=c_1^{1/2}4^{-4}(48N)^{-1}$, and $Z^{\mathrm{FE}}_k/Z_k'\to0$. Moreover,
for every $T>0$, no bound of Theorem~\textup{\ref{rU7:torus}(iv)(b)} holds for
the constants $Z_k'$: in the torus setting
\begin{gather*}
\sup_kL^{-4}\bigl(\log Z_{M_k}(\beta_k+h_k)-2\log Z_{M_k}(\beta_k)\\
+\log Z_{M_k}(\beta_k-h_k)\bigr)=\infty,\qquad
h_k=\frac{4TZ'_ka_k^4}N ,
\end{gather*}
and in the infinite-lattice setting $\sup_k\Xi^\infty_k(T)=\infty$ for the
constants $Z_k'$.

\textup{(iii)} If \textup{(AF$^\square$)} holds for constants $Z_k'$, then every
limit $A$ of Theorem~\textup{\ref{rU7:limit}} has $\mathrm EA(f)A(g)=0$ whenever
the supports of $f,g\in\mathcal R(X)$ are at positive distance, and there is
$\sigma^2\ge0$ with
\[
\mathrm EA(f)A(g)=\sigma^2\int_Xf(x)g(x)\,dx\qquad(f,g\in\mathcal R(X)).
\]

\textup{(iv)} Under the hypothesis of \textup{(iii)}, for $f,g\in\mathcal R(X)$
vanishing outside a compact subset of $P_{\mu,c}$,
$\mathrm EA(\theta_{\mu,c}f)A(g)=0$. In the Hilbert space obtained from the
positive form $(F,F')\mapsto\mathrm E[\overline{F(A(\theta_{\mu,c}f_j))}F'(A(
g_j))]$ of Theorem~\textup{\ref{rU7:limit}(iv)} (polynomials in fields
smeared in $P_{\mu,c}$, modulo null vectors), every single smeared field
$A(f)$ is a null vector, since its norm is $\mathrm EA(\theta_{\mu,c}f)A(f)
=0$: no one-field vector survives.
\end{theorem}

\begin{proof}
(i) $\min(h,1)\le1$.

(ii) For $|z|_\infty=2$ along a coordinate axis, $a_k|z|=2a_k\le r_0$ for large
$k$, and $|G^\square_k(z)|\le\|\Phi^\square\|_\infty^2\le(48N)^2$ because $0\le
\Phi^\square\le48N$; the lower half gives $Z_k'^2(48N)^2\ge c_1(2a_k)^{-8}(1+
\log\frac1{2a_k})^{-\gamma}$, which gives the bound on $Z_k'$ (with room to
spare in the constant). By Corollary~\ref{rU7:fe} or
Theorem~\ref{rU7:zfour}(ii), the hypothesis of Lemma~\ref{rU7:incompat} holds
for $Z_k=Z^{\mathrm{FE}}_k$ and $Q_*=[0,1)^4$ with $C=2e^{16}$ (on $\mathbb T_L$, with
the image of $[0,\min(1,L/2))^4$ and $C=2e^{L^4}$),
so $Z^{\mathrm{FE}}_k/Z_k'\to0$. If the displayed supremum were finite for some
$T$, then by Remark~\ref{rU7:torusremark} $K_T=\sup_k\Xi_{M_k,\beta_k,Z'_k}(T)
<\infty$, and by Theorem~\ref{rU7:torus}(i) with $t=\pm T/\|f\|_\infty$ and
$t^2y^2\le e^{ty}+e^{-ty}-2$, for bounded $f\ne0$,
\begin{align*}
\mu_k\bigl(A_{k,Z_k'}(f)^2\bigr)&\le\frac{2\|f\|_\infty^2}{T^2}\bigl(e^{K_T\|f\|
_{1,a_k}/\|f\|_\infty}-1\bigr)\\
&\le\frac{2K_Te^{K_TL^4}}{T^2}\|f\|_\infty\|f\|_{1,a_k};
\end{align*}
Lemma~\ref{rU7:incompat} with $Z_k=Z_k'$ would give $1\to0$. The
infinite-lattice case is the same with Theorem~\ref{rU7:zfour}(i) and
$\mathfrak m_k(f)\le2^4$ for $f$ vanishing outside $[0,1)^4$, by~\eqref{rU7:count}.

(iii) Let the supports of $f,g$ be at distance $\delta>0$. For large $k$, every
pair with $f(a_km_x)g(a_km_y)\ne0$ has $|y-x|_\infty\ge2$ and $a_k|y-x|\ge
\delta/2$, where $w_\gamma\le(\delta/2)^{-8}$; the upper half gives
\begin{align*}
&\bigl|\nu_k\bigl(A_k(f)A_k(g)\bigr)\bigr|\\
&\quad=\Bigl(\frac{Z^{\mathrm{FE}}_k}{Z_k'}
\Bigr)^2\Bigl|Z_k'^2\sum_{x,y}a_k^8f(a_km_x)g(a_km_y)G^\square_k(y-x)\Bigr|\\
&\quad\le\Bigl(\frac{Z^{\mathrm{FE}}_k}{Z_k'}\Bigr)^2\frac{2^8C_1}{\delta^8}\|f\|_{1,k}
\|g\|_{1,k}\to0
\end{align*}
by (ii), and the left side tends to $\mathrm EA(f)A(g)$ by
Theorem~\ref{rU7:limit}(i). For boxes $B,B'$ (products of half-open intervals)
with $B\cap B'=\emptyset$ let $B_\varepsilon\subset B$ be the box of the points
at $\ell^\infty$-distance at least $\varepsilon$ from the complement of $B$;
then $\mathrm EA(\mathbf 1_{B_\varepsilon})A(\mathbf1_{B'})=0$, and by the
second-moment bound of Theorem~\ref{rU7:limit}(ii) and Cauchy--Schwarz,
$|\mathrm EA(\mathbf1_{B\setminus B_\varepsilon})A(\mathbf1_{B'})|\le C(|B\setminus
B_\varepsilon|\,|B'|)^{1/2}\to0$; so $\mathrm EA(\mathbf1_B)A(\mathbf1_{B'})=0$.
Let $m(B)=\mathrm EA(\mathbf1_B)^2$. By linearity and these cross terms, $m$ is
additive over finite disjoint unions that are boxes; by
Theorem~\ref{rU7:limit}(iii) it is translation invariant. A cube $Q$ of side
$2^{-n}\ell$ tiles the cube of side $\ell$ by $2^{4n}$ translates, so
$m(Q)=|Q|\,\sigma^2$ with $\sigma^2=m(Q_0)/|Q_0|$ for a fixed cube $Q_0$
(on $\mathbb T_L$, $Q_0=\mathbb T_L$), and every box with corners in $b_0+2^{-n}\ell_0\Z^4$, $b_0$ a corner and
$\ell_0$ the side of $Q_0$, is a disjoint union of such cubes, so $m(B)=
\sigma^2|B|$ for these boxes. For a general box $B$ and such boxes $B_j$ with
$|B\triangle B_j|\to0$, $|m(B)^{1/2}-m(B_j)^{1/2}|\le\mathrm E(A(\mathbf1_B-
\mathbf1_{B_j})^2)^{1/2}\le(C|B\triangle B_j|)^{1/2}\to0$. Hence for step
functions on a common family of disjoint boxes, $\mathrm EA(f)A(g)=\sum_if_ig_i
\sigma^2|B_i|=\sigma^2\int fg$, and the identity extends to $\mathcal R(X)$ by the
approximation of Lemma~\ref{rU7:approx} and the $L^2$ bounds of Step~2 of the
proof of Theorem~\ref{rU7:limit}, since $\int fg$ is continuous under $L^1$
convergence with bounded sup norms.

(iv) $\theta_{\mu,c}f$ vanishes outside a compact subset of $\theta_{\mu,c}
(P_{\mu,c})$, which is at positive distance from that of $g$; apply (iii).
\end{proof}

\begin{remark}\label{rU7:afremark}
Theorem~\ref{rU7:afthm} makes the price of Theorem~\ref{rU7:limit} exact. Every
normalisation that controls the coinciding-point moments of the smeared
plaquette field with the two-norm bound $\|f\|_\infty\|f\|_{1,k}$ is
negligible against any normalisation for which the short-distance law holds;
under the full law the limit sees only a contact term, its covariance is
that of white noise, and its single smeared fields carry no vector of the
reflection-positive Hilbert space. This is the lattice counterpart of
Proposition~\ref{prop:limitEsing}, now with no hypothesis on the limit: the
renormalised field of (E) must be controlled at separated points only, as
in (M$_{\mathrm s}$). Nothing in (iii) says that the higher correlations of the
limit vanish at separated points, or that it is Gaussian; whether
$\sigma^2>0$ is the free-energy condition of Theorem~\ref{rU7:limit}(v).
\end{remark}
\subsubsection{Partition functions, sources, and what remains}
\label{rU7:balaban}

\emph{What is known.} Ba{\l}aban's renormalisation group analysis of
four-dimensional lattice gauge theories on a periodic
lattice~\cite{Bal84a,Bal85a,Bal85c,Bal87,Bal88,Bal89b} culminates, in the
last of these papers, in the statement that if the sequence of effective
coupling constants stays in an interval $(0,\gamma]$ with $\gamma$ small,
the effective densities produced by the renormalisation transformations
have the form and satisfy the bounds of his inductive scheme, from which
ultraviolet stability bounds of the same type as for superrenormalisable
models follow (we rely here on the review of~\cite{Bal89b} in zbMATH, not
on a reading of the original papers; no theorem of these papers is quoted
below, and their bounds enter only through hypothesis~(B), which is
assumed). Jaffe and Witten, describing Ba{\l}aban's analysis of the
three-dimensional lattice theory on a torus, call its estimates uniform in
the lattice spacing and note that ``these results need to be extended to the
study of expectations of gauge-invariant functions of the fields''; they
describe the four-dimensional extension~\cite{Bal87} as including an
analysis of asymptotic freedom and of large fields~\cite{JW06}. Magnen, Rivasseau and
S\'en\'eor~\cite{MRS93} construct the Schwinger functions of the pure
$\SU(2)$ gauge field in four dimensions with an infrared cut-off and no
ultraviolet cut-off, in a regularised axial gauge whose momentum cut-off
breaks gauge invariance and is compensated by counterterms, and verify the
Slavnov identities non-perturbatively; the infrared cut-off is not removed and
the Osterwalder--Schrader axioms are not established there. Gauge-invariant
smoothed observables via the Yang--Mills heat flow, proposed by Charalambous
and Gross (see~\cite{CC24}) and used in lattice gauge theory by
L\"uscher~\cite{Lus10}, are the basis of the state space of Cao and
Chatterjee for three-dimensional theories~\cite{CC24,CC23}, which defines
regularised Wilson loops under a Gaussian-free-field-like assumption on the
theory; Chevyrev~\cite{Che19} and Chandra, Chevyrev, Hairer and
Shen~\cite{CCHS22,CCHS24} construct state spaces of distributional connections
in two and three dimensions. None of these gives a bound on correlations of
gauge-invariant observables of the four-dimensional Wilson measure along a
lattice tuning. O'Carroll and Faria da Veiga~\cite{OCFdV20} give, by multiple
reflections, bounds on a normalised generating functional of a fixed number
$r$ of unsmeared gauge-invariant plaquette fields that are uniform in the
volume, the spacing and the coupling in a bounded range; since $r$ is
fixed, so these bounds do not control fields smeared over
$\asymp a^{-4}$ plaquettes.

\emph{From partition functions to correlations.} Theorem~\ref{rU7:torus}(iv)
and Remark~\ref{rU7:torusremark} are the exact interface: uniform exponential
moments of all smeared plaquette fields with normalisation $Z_k$ on
$\mathbb T_L$ are equivalent to the bound
\begin{equation}
\begin{split}
&\sup_kL^{-4}\bigl|\Delta^2_{h_k(T)}\log Z_{M_k}(\beta_k)\bigr|<\infty\quad
(T>0),\\
&h_k(T)=\frac{4TZ_ka_k^4}N ,
\end{split}
\tag{rU7.4}\label{rU7:delta}
\end{equation}
a statement about the partition function at three homogeneous couplings. A
stability bound of Ba{\l}aban's type feeds into~\eqref{rU7:delta} as follows.

\begin{corollary}[stability bounds at three couplings]\label{rU7:B}
Suppose that there are $K,E<\infty$ and functions $F_k$ on
$W_k=[\beta_k-a_k^2,\beta_k+a_k^2]$, $W_k\subset[0,\infty)$, such that
\begin{equation}
\begin{split}
&\bigl|\log Z_{M_k}(\beta')-F_k(\beta')\bigr|\le KL^4\quad(\beta'\in W_k),\\
&\bigl|F_k(\beta_k+H)-2F_k(\beta_k)+F_k(\beta_k-H)\bigr|\le EL^4\frac{H^2}{a_k^4}\\
&\qquad(|H|\le a_k^2).
\end{split}
\tag{B}
\end{equation}
Then $a_k^{-4}\max_\pm R_{M_k}(\beta_k;\pm H)\le EH^2/a_k^4+4K$ for $|H|\le a_k^2$,
$Z^{\mathrm{FE}}_k\ge N/(4(E+4K)a_k^2)$ when $a_k\le1$ and $E+4K\ge1$, and the fields
with $Z_k=N/(4a_k^2)$ have uniformly bounded exponential moments:
$\log\mu_k(e^{t\mathcal A_k(f)})\le(E+4K)\,|t|\,\|f\|_{1,a_k}$ for
$|t|\,\|f\|_\infty\le1$.
\end{corollary}

\begin{proof}
Both remainders are nonnegative and $M^4(R_M(\beta;H)+R_M(\beta;-H))=\Delta^2_H
\log Z_M(\beta)$, so, with $M_k^4=L^4a_k^{-4}$,
\[
a_k^{-4}\max_\pm R_{M_k}(\beta_k;\pm H)\le L^{-4}\Delta^2_H\log Z_{M_k}(\beta_k)
\le\frac{EH^2}{a_k^4}+4K .
\]
For $|H|\le a_k^2$ the right side is at most $E+4K$, and by convexity
$a_k^{-4}\max_\pm R(\pm sa_k^2)\le s(E+4K)$ for $s\in[0,1]$; so $h_{M_k,\beta_k}
\ge a_k^2/(E+4K)$, which gives the bound on $Z^{\mathrm{FE}}_k$ as
$a_k^2/(E+4K)\le1$. With $Z_k=N/(4a_k^2)$, $4TZ_ka_k^4/N=Ta_k^2$ and
$\Xi_{M_k,\beta_k,Z_k}(T)\le T(E+4K)$ for $T\le1$; Theorem~\ref{rU7:torus}(i)
gives the last claim.
\end{proof}

Ultraviolet stability bounds of the type described above would be
expected to give, at each fixed coupling, an estimate of the first kind
in~(B) with $F_k$ the extracted vacuum energy; we have not verified this
against the original papers. The uniformity of such an estimate over a
window of couplings of width $a_k^2$, and the second condition in~(B),
regularity of the extracted vacuum energy on that window, are further
inputs that we have not found stated in the descriptions of these papers
available to us. Corollary~\ref{rU7:B} is therefore a conditional statement
with (B) as its hypothesis. Granted~(B), the normalisation $a_k^{-2}$ is
reached, but no further:

\begin{corollary}[no partition-function route to separated correlations]
\label{rU7:nogo}
If the lower half of \textup{(AF$^\square$)} holds for the torus states with
constants $Z_k'$, then \eqref{rU7:delta} fails for $Z_k=Z_k'$ and every $T>0$,
while $h_k(T)\ge4cT(1+\log\frac1{a_k})^{-\gamma/2}/N$: no bound on homogeneous
partition functions at couplings $\beta_k$ and $\beta_k\pm h_k$, however
precise, controls the renormalised field of the short-distance law through the
criterion of Theorem~\ref{rU7:torus}(iv), although the coupling shifts
involved are bounded below by a negative power of $1+\log\frac1{a_k}$
only.
\end{corollary}

\begin{proof}
Theorem~\ref{rU7:afthm}(ii).
\end{proof}

\emph{The additional input.} What is needed is the partition function with an
inhomogeneous coupling. For bounded $f$ and $t\in\R$, by the definition of
$\mathcal A$ and $\sum_xf(am_x)\Phi^\square_x=\sum_p\psi_pw_p(f)$ with
$w_p(f)=\sum_{x:\,C_x\supset p}f(am_x)$,
\begin{align*}
\mu_{M,\beta}\bigl(e^{t\mathcal A(f)}\bigr)={}&e^{-tZa^4c\sum_xf(am_x)}\,
e^{tZa^4N\sum_pw_p(f)}\\
&\times\frac{Z_M[(\beta-tZa^4w_p(f)/N)_p]}{Z_M(\beta)} ,
\end{align*}
where
\[
Z_M[(\beta_p)_p]=\int\exp\Bigl(N\sum_p\beta_p\ReTr Q_p\Bigr)\,dQ .
\]
The mixed second derivative in $(s,u)$ at $0$ of
\[
\log Z_M\bigl[\bigl(\beta-sZa^4w_p(f)/N-uZa^4w_p(g)/N\bigr)_p\bigr]
\]
is the connected correlation $\mu(\mathcal A(f)\mathcal A(g))$; at the
normalisation of (AF$^\square$), by Theorem~\ref{rU7:afthm}(ii), the coupling
profile is a bump of height at least of order
$(1+\log\frac1{a_k})^{-\gamma/2}$ on the physical supports of $f$ and $g$. The input that would pass
from the renormalisation group to correlations of the smeared plaquette field
is therefore a stability bound for $\log Z_M$ with a coupling varying on
physical scales, with an error in the mixed second derivative that tends to
$0$ for separated $f,g$, rather than an error $O(L^4)$ in each partition
function. By Lemma~\ref{rU7:incompat} no bound of the two-norm form
$C\|f\|_\infty\|f\|_{1,k}$ on the variances holds at that normalisation, so
such an input can only be posed at separated supports, as (M$_{\mathrm s}$)
is for $n=2$.

\emph{Reconstruction in one direction.} The proved properties of the limits
give, in the infinite-lattice setting, the Hilbert space and the transfer
semigroup of one coordinate direction.

\begin{proposition}[one-direction reconstruction]\label{rU7:transfer}
Let $A$ be a limit of Theorem~\textup{\ref{rU7:limit}} in the
infinite-lattice setting, $\mu$ a direction, $\theta=\theta_{\mu,0}$,
$\tau_t=\tau_{te_\mu}$, and let $\mathfrak P_+$ be the set of random
variables $F=\Phi(A(f_1),\dots,A(f_n))$ with $\Phi:\R^n\to\C$ bounded and
continuous and $f_j\in\mathcal R$ vanishing outside a compact subset of
$P_{\mu,0}$; put $\theta F=\Phi(A(\theta f_1),\dots,A(\theta f_n))$ and
$\tau_tF=\Phi(A(\tau_tf_1),\dots,A(\tau_tf_n))$ for $t\ge0$. Then
$\langle F,F'\rangle=\mathrm E[\overline{\theta F}\,F']$ is a positive
semidefinite Hermitian form on $\mathfrak P_+$; let $\mathcal H_\mu$ be the
completion of $\mathfrak P_+$ modulo its null space, $[F]$ the class of $F$
and $\Omega=[1]$. There is a self-adjoint operator $H_\mu\ge0$ on
$\mathcal H_\mu$ with $[\tau_tF]=e^{-tH_\mu}[F]$ for $t\ge0$, $F\in\mathfrak
P_+$, and $H_\mu\Omega=0$.
\end{proposition}

\begin{proof}
Positivity is Theorem~\ref{rU7:limit}(iv) with $c=0$. Since $\theta\in B_4$,
Theorem~\ref{rU7:limit}(iii) shows that the families $(A(f_j),A(\theta f'_i))$
and $(A(\theta f_j),A(f'_i))$ have the same law, so
\[
\langle F',F\rangle=\overline{\langle F,F'\rangle},
\]
and the Cauchy--Schwarz inequality holds. For test functions
\[
\theta(\tau_tf)=\tau_{-t}(\theta f);
\]
translating the
family $(A(\tau_{-t}\theta f_j),A(\tau_tf_j))_j$ by $te_\mu$ and using
Theorem~\ref{rU7:limit}(iii),
\[
\langle\tau_tF,\tau_tF\rangle=\langle F,\tau_{2t}F\rangle\le\|F\|\,
\|\tau_{2t}F\| ,
\]
where $\|F\|=\langle F,F\rangle^{1/2}$. Applying this with $t$ replaced by $2t,4t,\dots$,
\[
\|\tau_tF\|\le\|F\|^{1-2^{-m}}\|\tau_{2^mt}F\|^{2^{-m}}\qquad(m\ge1),
\]
and
$\|\tau_sF\|\le\sup|\Phi|$ for every $s$; letting $m\to\infty$ gives
$\|\tau_tF\|\le\|F\|$. Hence $T_t[F]=[\tau_tF]$ is well defined (null
vectors go to null vectors) and extends to a contraction of $\mathcal H_\mu$,
with $T_0=1$, $T_sT_t=T_{s+t}$ and $T_t\Omega=\Omega$. Translating
$(A(\theta g_i),A(\tau_tf_j))$ by $-te_\mu$ gives $\langle G,\tau_tF\rangle=
\langle\tau_tG,F\rangle$, so $T_t$ is self-adjoint, and $T_t=T_{t/2}^2\ge0$.
For strong continuity, by Cauchy--Schwarz in $L^2(\mathrm P)$ and the
equality of the laws of $(A(\theta\tau_tf_j),A(\theta f_j))_j$ and $(A(\tau_t
f_j),A(f_j))_j$,
\[
\|T_t[F]-[F]\|^2\le\mathrm E|\tau_tF-F|^2 ,
\]
which tends to $0$ as $t\to0$
because $A(\tau_tf_j)\to A(f_j)$ in $L^2(\mathrm P)$
(Theorem~\ref{rU7:limit}(iii)) and $\Phi$ is bounded and continuous. Let $T=T_1$; then $0\le T\le1$. If
$T\psi=0$ then, by induction on $m$, $\|T_{2^{-m-1}}\psi\|^2=\langle\psi,
T_{2^{-m}}\psi\rangle=0$, so $T_{2^{-m}}\psi=0$ for all $m$, hence $\psi=\lim_mT_{2^{-m}}\psi=0$. So $T$ is injective and
$H_\mu=-\log T$, defined by the spectral theorem~\cite[Sec.~VII.2]{RS80},
is self-adjoint and $H_\mu\ge0$. Since $T_{2^{-m}}\ge0$ and
$(T_{2^{-m}})^{2^m}=T$, $T_{2^{-m}}$ is the positive $2^m$-th root
$e^{-2^{-m}H_\mu}$ of $T$, so $T_t=e^{-tH_\mu}$ for dyadic $t\ge0$, and for all
$t\ge0$ by the strong continuity of both sides; $T_t\Omega=\Omega$ gives
$H_\mu\Omega=0$.
\end{proof}

On the torus the same form gives a Hilbert space and the vacuum $\Omega$, but
the positive region $P_{\mu,c}$ is bounded and $\tau_t$ does not map it into
itself for every $t>0$, so no semigroup is claimed there. By
Theorem~\ref{rU7:afthm}(iv), under (AF$^\square$) every single smeared field is
a null vector of $\mathcal H_\mu$.

\emph{What remains for (E), reconstruction, and nontriviality.} For the limits
of Theorem~\ref{rU7:limit} the following are proved: existence along
subsequences for every coupling sequence, translation and $B_4$ invariance,
reflection positivity in every coordinate direction, centring, moment
bounds $|\mathrm E\prod_jA(f_j)|\le2\,n!\,e^n\prod_j(\|f_j\|_\infty+\|f_j\|
_{L^1})$, which are of the type of the linear growth condition
of~\cite{OS75} for product test functions with a norm of order zero (we do
not show that the moments extend to tempered distributions), and, on
$\R^4$, the reconstruction of Proposition~\ref{rU7:transfer} in each
coordinate direction. Not proved, and needed:
(a) $\mathrm O(4)$ invariance; the set of $R\in\mathrm O(4)$ leaving the law of
the limit invariant is a subgroup containing $B_4$, closed because $f\mapsto
A(f)$ is continuous into $L^2(\mathrm P)$ under $L^1$ convergence with bounded
sup norms and supports in a fixed cube (Step~2 of the proof of
Theorem~\ref{rU7:limit}), so by the group-theoretic argument of
Theorem~\ref{thm:limitEsep}(vii) one rotation of infinite order would
suffice; (b) independence of the subsequence; (c) clustering; (d) a
normalisation that sees separated points: by Theorem~\ref{rU7:afthm}, under
(AF$^\square$) the limits of Theorem~\ref{rU7:limit} have white-noise
covariance and no one-field vectors, so a nontrivial two-point function at
separated points requires $Z_k/Z_k^{\mathrm{FE}}\to\infty$ and bounds posed at
separated supports only, as in the inputs \ref{rE:I2}--\ref{rE:I5} of
Sec.~\ref{ssec:rEmissing}; (e) for the mass gap, the averaged condition
\ref{rE:I1} along the tuning; (f) for nontriviality in the sense of a non-free
theory, beyond $S_2\ne0$ at separated points, a nonvanishing connected
four-point function, since a generalised free field satisfies all the other
requirements. None of (a)--(f) is proved here, and neither the mass gap, nor
(D), nor (E) is claimed.

\emph{Bounded observables.} Wilson loops of physical size, with or without
gradient flow, are bounded by $N$ in absolute value, so their joint laws
along any tuning are tight; for them the content of (E) lies entirely in
renormalisation and nontriviality, which for spatially flowed packets is
the ratio condition of Proposition~\ref{rE:flow}.
Theorem~\ref{rU7:limit} is the corresponding statement for an unbounded
field, where tightness itself has content: it requires the free-energy
normalisation, and Theorem~\ref{rU7:afthm} shows that under (AF$^\square$)
this normalisation is the obstruction.
